\documentclass[letterpaper,12pt]{article}
\usepackage[T1]{fontenc}
\usepackage[utf8]{inputenc}
\usepackage[spanish]{babel}
\usepackage{framed}
\usepackage{amsfonts,setspace}
\usepackage{amsmath}
\usepackage{amssymb, amsmath, amsthm}
\usepackage{comment}
\usepackage[most]{tcolorbox}
\usepackage{amssymb}
\usepackage{dsfont}
\usepackage{anysize}
\usepackage{pdfpages}
\usepackage{tikz}
\usepackage{hyperref}
\usepackage{multicol}
\usepackage{enumerate}
\usepackage{graphicx}
\usepackage{float}
\usepackage[dvipsnames]{xcolor}
\usepackage[left=1.5cm,top=0.9cm,right=1.5cm, bottom=1.4cm]{geometry}
\usepackage{fancyhdr}
\pagestyle{fancy}
\definecolor{blue(pigment)}{rgb}{0.2, 0.2, 0.6}
% Page
\oddsidemargin -0.4in
\textwidth 7in
\topmargin -0.6in
\textheight 9.1in
\parindent 0em
\parskip 1ex

%Teoremas, Lemas, etc.
\theoremstyle{plain}
\newtheorem{teo}{Teorema}
\newtheorem{lem}{Lemma}
\newtheorem{prop}{Proposici\'on}
\newtheorem{cor}{Corolario}
\theoremstyle{definition}
\newtheorem{defi}{Definici\'on}
\newtheorem{eje}{Ejemplo}
\newtheorem{obs}{Observaci\'on}
% fin macros

% Comandos más usados
% Conjuntos
\newcommand{\N}{\mathbb{N}}
\newcommand{\Z}{\mathbb{Z}}
\newcommand{\Q}{\mathbb{Q}}
\newcommand{\R}{\mathbb{R}}
\newcommand{\C}{\mathbb{C}}
\newcommand{\RR}{\mathbb{R} \times \mathbb{R}}
\newcommand{\fl}{\operatorname{fl}}
% Flechas
\newcommand{\ssi}{\Longleftrightarrow}
\newcommand{\imp}{\Longrightarrow}
\newcommand{\pmi}{\Longleftarrow}
% Trigonométricas
\newcommand{\sen}{\text{sen}}

\usepackage{versions}
\includeversion{solution}

\let\epsilon\varepsilon

\begin{document}
\def\Pauta{1}
%%%%%ESCRIBIR 1 si quieren ver la Prueba con Pauta%%%%
%%%%%ESCRIBIR 0 si quieren ver sólo la Ayudantía%%%%
	
	\fancyhead[L]{\itshape{Escuela de Ingeniería}}
	\fancyhead[R]{\itshape{Universidad de O'Higgins}}
	
	\begin{minipage}{11.5 cm}
         \vspace{-6mm}
		\begin{flushleft}
			\vspace{1mm}
			\textbf{ING2602-1  Métodos Numéricos}\\
			\textbf{Profesor:} Gonzalo Flores  \\
			\textbf{Ayudante:}  Juan I. Riquelme \\
		\end{flushleft}
	\end{minipage}
	
	% \begin{picture}(2,3)
	% 	\put(430,9.5){\includegraphics[scale=0.22]{UOH.png}}
	% \end{picture}
	\vspace{-6mm}
	\begin{center}
		\LARGE \bf{Ayudantía 3: Punto Flotante} %Teorema Fundamental del Cálculo
	\end{center}
\begin{enumerate}[\textbf{P\arabic{enumi}.}]

#Preguntas
\item Ya que en el computador solamente se pueden representar una cantidad finita de números, utilizamos conjuntos finitos $F$, llamados punto flotante. Un conjunto de punto flotante $F$ se define mediante una base $\beta \in \mathbb{N}$, en un intervalo $[L,U]$ del exponente $E \in \mathbb{Z}$ y una precisión $p \in \mathbb{N}$.

Consideramos un \textit{conjunto de punto flotante normalizado} con los siguientes valores:
\[
    \beta = 3,
    \qquad
    L = -1,
    \qquad
    U = 2,
    \qquad
    p = 2.
\]

\end{enumerate}


%Resumen

\begin{tcolorbox}[
    title = Resumen,
    colframe=cyan!80!black,
    colback=cyan!5!white,
    boxrule=0.8pt,
    arc=4pt,
    auto outer arc
]

\begin{multicols}{2}

\defi[Sistema de punto flotante.]
Un sistema de punto flotante queda determinado por
\[
    F(\beta,p,L,U),
\]
donde $\beta$ es la base, $p$ la cantidad de cifras significativas y
$L\leq E\leq U$ es el rango permitido para el exponente.

Un número normalizado distinto de cero tiene la forma
\[
    x=\pm(d_0.d_1\ldots d_{p-1})_{\beta}\,\beta^E,
\]
donde
\[
    d_0\neq 0,
    \qquad
    L\leq E\leq U.
\]
Por lo tanto, \textbf{normalizar} significa escribir el número de manera
que la primera cifra de la mantisa sea distinta de cero. Por ejemplo,
\[
    202_3=2.02_3\times 3^2.
\]
\end{multicols}

\end{tcolorbox}

\end{document}
